\documentclass[11pt]{article}
\usepackage{../shared/luxcover}
\usepackage{amsmath,amssymb}
\usepackage{hyperref}
\usepackage{booktabs}
\usepackage{enumitem}
\usepackage{xcolor}

\title{Lux Identity}
\author{Zach Kelling\\\textit{Lux Industries, Inc.}\\\texttt{research@lux.network}}
\date{October 2023}

\begin{document}

\luxcoverpage

\begin{abstract}
Zero Knowledge Sovereign Identity

Lux Identity Provides a solution to automate the KYC/AML process through leveraging homomorphic encryption to secure private data into a zero knowledge file system. This tool gives Web3 application engineers an ability to build compliant and regulated systems with easy to install API integrations.
\end{abstract}

\tableofcontents
\newpage


Lux Identity Provides a solution to automate the KYC/AML process through leveraging homomorphic encryption to secure private data into a zero knowledge file system. This tool gives Web3 application engineers an ability to build compliant and regulated systems with easy to install API integrations.


\subsection{Self Sovereign Identity (SSID)}

A self-sovereign identity (SSID) is an online identity that a person can control themselves. This means that the individual has the power to determine what information is shared about them and with whom. The individual also has the ability to revoke access to this information at any time. A self-sovereign identity is different from a traditional online identity in that it does not rely on a central authority to manage and control the data. This gives the individual a great deal of control over their own identity and data. There are many benefits to having a self-sovereign identity.


\subsection{Benefits of SSID VS Other Identity Systems}

\subsubsection{Power to control Your Own Data}

One of the most significant benefits is that SSID gives individuals the power to control their own data. With traditional online identities, individuals have very little control over how their data is used and shared and they oftentimes have to share a great deal of personal information in order to use those applications. This can lead to situations where personal data is mishandled or even stolen. With a self-sovereign identity, individuals can choose what data they share and with whom they share it. This gives them a much greater degree of control over their online identity and helps to protect their personal data.


\subsubsection{Protect Community from Bad Actors}

Another benefit of self-sovereign identity is that it helps to protect the community. Through Lux Identity, companies and financial institutions can blacklist people who want to harm the community by blocking users who have been convicted of theft, fraud, money laundering or other crimes that conflict with an activity they wish to perform on the network.


\subsection{Lux Identity Is Both Trustless and Trusted}

A fully trustless economy has many inefficiencies --- you cannot build long-term relationships between people if there is no basis of trust.


We create an environment of trust without revealing unnecessary information. Additionally, Lux Identity allows for:


Identity verification --- making sure you are a legal citizen, from an accepted country, age, etc.


Sybil resilience --- one user cannot have multiple accounts, so one person for one vote, one person for one airdrop. Democracy.


Blacklisting bad actors --- filtering against public sanctions lists (e.g. US's SDN)


Compliance --- if a verified user does something illegal, their identity can be revealed and reported to the government


Privacy --- Lux Identity keeps your personal information confidential through leveraging ZChain as a Zero Knowledge File System.


\section{Benefits For Customers Using Lux Identity}

Zero-Knowledge means that you don't have to reveal any personal or identifiable information about yourself to anyone above and beyond what you're being verified for, e.g. verification that you're over 18 without revealing your birthday.


Decentralized means that you are in control of your data and no third party will have access to your personal information.


All data is stored as a reduction and at no point in time is viewable in plain text form using Zero Knowledge Proofs


\subsection{Zero Knowledge Proofs}

The zero-knowledge identity proof technique offers a way of verifying or providing proof of our identity whereby one party proves to know a particular piece of information without revealing other private information. One example of how Lux Identity uses the zero-knowledge proof protocol includes submitting proof of identity without disclosing your address or demonstrating that your bank account is sufficient for a particular transaction without revealing its balance.


A zero-knowledge identity proof is a term used to refer to an authentication scheme where one party proves to the other to have a particular piece of knowledge that proves ownership of the identity. The prover verifies the required information without further disclosing any additional sensitive or personal information. This ensures that you maintain ownership of your sensitive private data.


Zero-knowledge proof (ZKP) alerts the verifier that the prover has the required information to confirm his identity. The method was introduced during the 80s by MIT researchers and is used to further enhance blockchain functionality.


\section{Digital Signal Processing}

Biometric identification technology has become increasingly common in our daily lives as the demand for information security and security legislation has grown around the world. Due to its potential to overcome a number of fundamental constraints of unimodal biometric systems, multimodal biometrics technology has acquired interest and popularity in this respect. Lux Identity uses a new multimodal biometric human identification system based on a deep learning algorithm for identifying humans using biometric modalities of iris, face, and analog voice recognition.


Analog and digital signals are the two main forms of signals. A digital signal is a stream of 0s and 1s, whereas an analog signal is an uninterrupted, continuous stream of information.


DSP systems, which sample and digitize analog signals into a stream of discrete digital values using an Analog-to-Digital Converter, are one of the most significant components of biometric systems (ADC).


The DSP architecture is built to support complex mathematical algorithms that involve a significant amount of multiplication and addition.


It can also enhance the resolution of the captured image or audio file with the use of two-dimensional Fast Fourier Transforms (FFT) and finite IR filters.


\subsection{Deep Learning Features: Convolutional Neural Network for Anti-Spoofing}

Face spoofing detection is one of the most well-studied problems in computer vision. Face recognition has become a widely adopted technique in biometric authentication systems. In facial recognition based authentication techniques, the system first recognizes the person to verify the legitimacy of the user before granting access to the system resources. The system must be able to determine the liveness of the person in front of the camera, for example, by recognizing the face and denying the types of face presentation attacks related to photographs, videos and the 3D mask of the targeted person. Attackers try to directly or indirectly, masquerade the biometric system as another person by forging biometric traits in attempting to gain unauthorized access.


In response to these types of attacks, Deep learning and convolutional neural networks (CNN) are additional solutions that can help with anti spoofing.


Convolutional neural networks are a type of neural network that are used to learn patterns from data. They are similar to traditional neural networks, but they have a few key differences. Convolutional neural networks are made up of neurons that have a local connectivity pattern. This means that they are only connected to a small subset of the input data. They also have a weight sharing mechanism. This means that the weights of the neurons are shared across the input data. Lastly, convolutional neural networks use a pooling layer. This layer helps to reduce the dimensionality of the data. These characteristics make convolutional neural networks well suited for tasks such as image recognition and claSSIDfication.


We started thinking of anti-spoofing as a binary claSSIDfication problem when exploring technologies. We train the CNN to recognize which are live video and analog audio signals and which are spoofed.


The Lux biometrics authentication system is built around convolutional neural networks (CNNs), which extract features and use a softmax claSSIDfier to identify images. Three CNN models, one for the iris, one for the face, and one for Analog voice recognition, were integrated to create the Lux Identity system. Image augmentation and dropout techniques were employed to avoid overfitting while machine learning algorithms are implemented to stop identity spoofing.


Different fusion approaches were used to fuse the CNN models in order to understand the effects of fusion approaches on recognition performance; thus, feature and score level fusion approaches were used. Several tests using the SDUMLA-HMT dataset, a multimodal biometrics dataset, were used to empirically evaluate the proposed system's performance. The results showed that in biometric identification systems, employing three biometric qualities yielded better outcomes than utilizing two or one biometric traits. The results also revealed that our method surpassed other state-of-the-art methods, with an accuracy of 99.39 percent.


\section{Biometric DSP}

To provide Self sovereign identities, Lux identity leverages an open source DSP solution called BOB. Bob's primary data structure is a multi-dimensional array. Arrays are a good representation for many different types of digital signals in signal proceSSIDng and machine learning, such as images, audio data, and extracted features.


\section{Feature Extraction}

Lux uses feature extraction for the process of dimensionality reduction by which an initial set of raw data is reduced to more manageable groups for proceSSIDng. A characteristic of these large data sets is a large number of variables that require a lot of computing resources to process. Feature extraction is the name for methods that select and /or combine variables into features, effectively reducing the amount of data that must be processed, while still accurately and completely describing the original data set.


\section{Decentralized Infrastructure}

Lux Identity operates as a subnet of the Lux ZChain. ZChain is a zero knowledge privacy preserving subnetwork of Lux Network that utilizes zero knowledge proofs to allow for private transactions. Zero knowledge proof systems allow for a person to prove that they know something without revealing what that something is.


In the context of ZChain, this means that a person can prove that they own a certain asset without revealing what that asset is. This allows for private transactions to take place on the blockchain. ZChain is the first privacy preserving bridge that is compatible with Ethereum. This means that it can be used to create decentralized applications that require private transactions.


\subsection{Zero Knowledge Privacy Bridge}

Zero knowledge bridges bring privacy to DeFi by ensuring that data remains confidential between parties. zero knowledge circuits allow for the verifiable computation of encrypted data, while zkVM provides a secure execution environment for smart contracts.


ZChain is a decentralized network that enables the private exchange of data and value. By using ZChain, DeFi platforms can keep data confidential, preventing fraud and protecting users' privacy.


\section{Validator Fees}

A transaction fee is required for each vote as nodes process the Digital Signal Processing and verify video and audio signals against biometric records in the DSP database. These fees are charged to Lux Users as they use Lux identity to authenticate their identity to vote, make purchases, withdraw money from Lux ATMs, verify identity for access to accounts, documents, ect.


Each identity verification process is designed to operate using a low cost minimal fee that is equal to fractions of a percentage on each transaction.  With projected costs of less then 1 cent per transaction, Lux Identity is designed to be a lightweight security addition for applications looking to eliminate fraud from their communities.


\end{document}
